\section{Regularization from finite relative entropy}
\label{l:finite-entropy:section}

This argument restores a dyadic mass bound after every block.  The
input to its refinement procedure is any posterior of finite relative
entropy in its current cell.  The procedure may take an unbounded
number of steps, but its expected number of attempts and expected
increase in depth are finite.  These bounds pay for the complete
refinement label, including every decision to remain outside the
selected descendants.  Finite relative entropy supplies this expected
termination bound even when no uniform bound on the posterior density
is available.

Use the half-open cube \(Q_0\), uniform probability \(\mu_0\), and
graph map \(\phi\) of Section~\ref{l:dyadic}.  In this section write
\(Q_{\mathrm{ho}}=Q_0\), \(\mu=\mu_0\), \(s=\phi\), and
\(n=d-1\).  The dyadic cells remain half-open.
For a cell \(C\), let \(\mu_C\) be uniform probability on \(C\).
A descendant of relative depth \(h\) has \(\mu_C\)-mass \(2^{-nh}\).

For a probability \(P\ll\mu_C\) of finite relative entropy in a
depth-\(j\) cell, the change of reference measure gives
\begin{equation}\label{l:finite-entropy:cell-divergence}
 D(P\Vert\mu)=D(P\Vert\mu_C)+nj\log2,
\end{equation}
because \(d\mu_C/d\mu=2^{nj}\) on \(C\). The depth term is
the information certified by confinement to the cell; the remaining
divergence measures concentration within it. This identity holds even
when \(P\) is supported on a proper subset of \(C\), so complement
decisions remain represented in the current-cell divergence. The
refinement below bounds the expected remaining divergence and expected
depth increase by a single budget: the incoming divergence plus an
absolute constant. That budget also includes decisions that keep the
current cell.

\begin{definition}\label{l:finite-entropy:regular}
Put \(\alpha=1/2\).  A probability \(P\) supported on a cell \(C\)
is regular on \(C\) if every descendant \(E\) of relative depth
\(h\ge1\) satisfies
\[
 P(E)\le2\,2^{-\alpha nh}.
\]
\end{definition}

\subsection{The refinement and its termination}

\begin{lemma}[Finite expected regularization]
\label{l:finite-entropy:regularization}
Suppose \(P\ll\mu_C\) and \(K_0=D(P\Vert\mu_C)<\infty\).
For all sufficiently large \(n\), a countable label determined by
\(W\sim P\) specifies a descendant \(C'\) and a conditional law
\(P'\) regular on \(C'\).  The label is the whole terminating
history in the construction below, including every complement
decision; \(P'\) is conditioned on this label and need not be
\(P(\,\cdot\mid C')\).

Let \(H\) be the increase in cell depth and \(A_*\) the number of
attempts.  Then
\[
 \mathbb EH<\infty,\qquad
 \mathbb EA_*\le2(1+\mathbb EH)<\infty,
\]
and, with \(b=(1-\alpha)(\log2)/2\),
\begin{equation}\label{l:finite-entropy:budget}
 \mathbb E D(P'\Vert\mu_{C'})+bn\,\mathbb EH\le K_0+2/e.
\end{equation}
The hypothesis imposes no boundedness on \(dP/d\mu_C\).
\end{lemma}

\begin{proof}
At an attempt, consider the current conditional law in its current
cell.  A strict descendant \(E\) of relative depth \(h\ge1\) is
heavy if its current mass \(m_E\) exceeds \(2^{-\alpha nh}\).
Select the heavy descendants with no heavy strict ancestor of positive
relative depth.  Every heavy cell has such a first heavy ancestor,
because its chain of ancestors back to the current cell is finite.
The selected cells are a countable disjoint family.

If \(W\) is in one of these cells, reveal which cell, condition on
it, replace the current cell by it, and begin another attempt.
Otherwise reveal that \(W\) is in the complement, whose current
probability is \(q\), condition on that complement, and keep the
current cell.  Stop if \(q\ge1/2\); if \(0<q<1/2\), begin another
attempt.  Zero-probability branches are omitted.  These operations
are Borel and the label records their finite sequence if the procedure
stops.

At a stop, a descendant that was heavy immediately before the last
complement decision lies in a selected cell, so it has zero posterior
mass.  Every other descendant of depth \(h\) had mass at most
\(2^{-\alpha nh}\), and conditioning on a complement of probability
at least \(1/2\) increases this by at most two.  The resulting law
is therefore regular.

It remains to prove that a stop occurs and that its cost satisfies
\eqref{l:finite-entropy:budget}.  Let \(f\) be the density of the
current law relative to the uniform probability on the current cell.
On a heavy transition of depth \(h\), its value at the sampled point
changes to
\[
 f(W)\frac{2^{-nh}}{m_E}.
\]
Its logarithm therefore decreases by at least
\((1-\alpha)nh\log2\).  On a complement transition it increases
by \(\log(1/q)\).  These are exact density updates on every
positive-probability branch, with all previous complement exclusions
retained in the conditional law.

Truncate after \(N\ge1\) attempts, freezing a path once it has stopped.
Let \(H_N\) be accumulated heavy depth, \(J_N\) the sum of complement
logarithmic increases, and \(A_N\) the number of attempts actually made.
If \(f_N\) is the resulting conditional density, the pathwise updates
give
\begin{equation}\label{l:finite-entropy:path}
 \log f_N(W)+(1-\alpha)n(\log2)H_N
 \le\log f_0(W)+J_N.
\end{equation}
For an active attempt the conditional expected complement increase,
counted as zero on a heavy branch, is \(q\log(1/q)\le1/e\).
Thus
\[
 \mathbb EJ_N\le \mathbb EA_N/e\le N/e.
\]

We justify taking expectations of the logarithms before using
\eqref{l:finite-entropy:path}.  For any density \(f\) relative to a
probability measure \(\lambda\),
\[
 \int_{\{f<1\}} f\log(1/f)\,d\lambda\le1/e,
\]
because \(x\log(1/x)\le1/e\) for \(0<x<1\).
This bounds the expected negative log part at the sampled point,
conditionally on each positive-probability history.  Finite \(K_0\)
then makes the positive part of \(\log f_0(W)\) integrable.
The path inequality bounds the positive part of \(\log f_N(W)\)
by \((\log f_0(W))^++J_N\).  Consequently
\(\mathbb E\log f_N(W)\) is well defined and is the expected
conditional relative entropy, hence is nonnegative.  Taking
expectations in \eqref{l:finite-entropy:path} first shows
\(\mathbb EH_N<\infty\).

An attempt after the first follows a heavy transition or a continuing
complement.  The number of heavy transitions is at most \(H_N\),
because each adds at least one unit of depth.  At each active attempt,
the conditional probability of a continuing complement is at most
\(1/2\).  Counting such predecessors and taking expectations gives
\begin{equation}\label{l:finite-entropy:attempts}
 \mathbb EA_N\le1+\mathbb EH_N+\tfrac12\mathbb EA_N,
 \qquad
 \mathbb EA_N\le2(1+\mathbb EH_N).
\end{equation}
Substituting \(\mathbb EJ_N\le\mathbb EA_N/e\) into the expected
path inequality now yields
\begin{equation}\label{l:finite-entropy:truncated}
 \mathbb E\log f_N(W)
 +\bigl((1-\alpha)n\log2-2/e\bigr)\mathbb EH_N
 \le K_0+2/e.
\end{equation}
For sufficiently large \(n\), the coefficient is at least
\(bn=(1-\alpha)n(\log2)/2\).  Hence the expectations of \(H_N\)
and, by \eqref{l:finite-entropy:attempts}, \(A_N\), are bounded
uniformly in \(N\).

Both sequences increase on every path.  Monotone convergence gives
finite expectations for their limits \(H,A_*\), and preserves
\(\mathbb EA_*\le2(1+\mathbb EH)\).  A path that never stops
has \(A_N\to\infty\), so nontermination has probability zero.
The possible terminating histories form a countable set of finite
sequences.

The limiting complement sum \(J\) also satisfies, by monotone
convergence,
\[
 \mathbb EJ\le\mathbb EA_*/e
 \le(2/e)(1+\mathbb EH)<\infty.
\]
At termination the positive log part of the final density is
integrable by the path inequality, and its negative log part is
integrable by the conditional bound above.  We may therefore take
expectations directly in the terminating path inequality.  This gives
\[
 \mathbb E D(P'\Vert\mu_{C'})
 +(1-\alpha)n(\log2)\mathbb EH
 \le K_0+(2/e)(1+\mathbb EH),
\]
which implies \eqref{l:finite-entropy:budget}.  In particular, no
passage of a signed logarithm through a limit is being used.
\end{proof}

\subsection{Local mass and exact projection densities}

We next bound the information from a block that starts with a regular
law.  The projection estimate is related to the independent-draw
Gaussian projection method of Sharan, Sidford, and Valiant
\cite[Section~7, Lemma~12]{SSV2019}.  The estimate below is proved
for exact multirow observations and for finite restrictions with their
density multiplier visible.

\begin{lemma}[Ambient ball bound]\label{l:finite-entropy:ball}
Suppose \(d\ge3\) and \(P\) is regular on a depth-\(j\) cell.
Let \(z\) be the graph image of its center and set \(R=2\,2^{-j}\).
Then the image of \(P\) is supported in \(B(z,R)\).  With
\(\beta=1/3\), there is an absolute \(C_0\ge1\) such that
\begin{equation}\label{l:finite-entropy:ball-bound}
 P\{w:\|s(w)-x\|\le r\}
 \le C_0^d(r/R)^{\beta d}
 \quad(x\in\mathbb R^d,\ 0<r\le R).
\end{equation}
\end{lemma}

\begin{proof}
The cell diameter is \(2^{-j}\); the \(2\)-Lipschitz property gives
the stated support radius.  For \(r\le R/2\), choose \(h\ge1\)
with \(r\le R2^{-h}<2r\).  The depth-\((j+h)\) cells have side
between \(r/(2\sqrt n)\) and \(r/\sqrt n\), and diameter at most
\(r\).  A ball test for the graph image restricts \(w\) to a ball
of radius \(r\) in its first \(n\) coordinates.  Every cell meeting
that ball lies in the concentric ball of radius \(2r\).  Comparing
disjoint cell volumes and using the unit-ball estimate from
Lemma~\ref{ld:cube-geometry} gives at most
\[
 v_n(2r)^n(2\sqrt n/r)^n\le(4\sqrt{2\pi e})^n
\]
such cells.  Each has \(P\)-mass at most \(2\,2^{-\alpha nh}\).
Since \(2^{-h}<2r/R\) and
\(\alpha n=(d-1)/2\ge d/3=\beta d\), all factors are bounded
by \(C_0^d(r/R)^{\beta d}\) for an absolute \(C_0\).
For \(R/2<r\le R\), enlarge \(C_0\) and use total mass one.
\end{proof}

\begin{lemma}[Projection with a density multiplier]
\label{l:finite-entropy:projection}
Fix \(\beta=1/3\), put \(\gamma=\beta/2=1/6\), and let
\(p=\lfloor\beta d/2\rfloor\ge2\).  Let \(1\le k\le\gamma d/2\),
and let \(\rho\) be the law of a \(k\)-by-\(d\) standard Gaussian
matrix.  Suppose that \(\nu\) is a finite positive Borel measure
of mass at most one, supported in \(B(z,R)\) for \(R>0\), and
that for some \(a\ge0\),
\[
 \nu(B(x,r))\le e^{ad}C_0^d(r/R)^{\beta d}
 \quad(x\in\mathbb R^d,\ 0<r\le R).
\]
The image of \(\rho(dA)\nu(ds)\) under \((A,s)\mapsto(A,As)\)
has a nonnegative density \(g\) relative to \(\rho(dA)\,dy\) with
\begin{equation}\label{l:finite-entropy:projection-moment}
 \int g(A,y)^p\,d\rho(A)\,dy
 \le\left[2R^{-k}e^{(a+C_1)k/\gamma}\right]^{p-1}.
\end{equation}
Here \(C_1\ge0\) depends only on \(C_0\), uniformly over all the
other parameters in the stated ranges.
\end{lemma}

\begin{proof}
Let \(F\) be an affine subspace of dimension \(m\le p-2\).
For \(0<r\le R\), a maximal \(r\)-separated set in the projection
of \(B(z,R)\) onto \(F\) has at most \((5R/r)^m\) points, by
packing disjoint radius-\(r/2\) balls.  Ambient balls of radius
\(2r\) at these centers cover the part of \(B(z,R)\) at distance
at most \(r\) from \(F\).  If \(2r>R\), the trivial mass bound
can replace the local bound at that radius, whose right side would
already be at least one.  Since \(\beta d-m\ge\gamma d\),
absorbing the covering constants gives
\begin{equation}\label{l:finite-entropy:tube}
 \nu\{s:\operatorname{dist}(s,F)\le r\}
 \le\min\{1,e^{(a+C_1)d}(r/R)^{\gamma d}\}.
\end{equation}
In particular \(F\) has zero \(\nu\)-mass.
Set \(r_0=R e^{-(a+C_1)/\gamma}\).  Layer-cake integration,
using total mass at most one above the cutoff \(r_0\), gives
\begin{equation}\label{l:finite-entropy:inverse}
 \begin{split}
 \int\operatorname{dist}(s,F)^{-k}\,d\nu(s)
 &\le r_0^{-k}
   +k r_0^{-\gamma d}\int_0^{r_0}r^{\gamma d-k-1}\,dr\\
 &=r_0^{-k}\left(1+\frac{k}{\gamma d-k}\right)
 \le2R^{-k}e^{(a+C_1)k/\gamma}=:B.
 \end{split}
\end{equation}
The last inequality uses \(k\le\gamma d/2\).

Choose a continuous nonnegative probability density \(\varphi_\eta\)
supported in \(B(0,\eta)\subseteq\mathbb R^k\), and define
\[
 g_\eta(A,y)=\int\varphi_\eta(y-As)\,d\nu(s).
\]
For \(\nu^{\otimes p}\)-almost every tuple \(s_1,\ldots,s_p\),
the differences \(s_i-s_1\), \(2\le i\le p\), are linearly
independent: at the \(i\)th step, the preceding affine span has
dimension at most \(i-2\le p-2\), so it is null by
\eqref{l:finite-entropy:tube}.  For such a tuple put
\(D=[s_2-s_1,\ldots,s_p-s_1]\).  The joint Gaussian vector of
its images under \(A\) has density at most
\[
 (2\pi)^{-k(p-1)/2}\det(D^{\mathsf T}D)^{-k/2}
 \le\prod_{i=2}^p
 \operatorname{dist}(s_i,\operatorname{aff}(s_1,\ldots,s_{i-1}))^{-k}.
\]
The determinant factorization is Gram--Schmidt.

Expand \(g_\eta^p\), shift the label integration variable by
\(As_1\), and integrate over the Gaussian difference images.
Bounding their density by the displayed supremum leaves a product of
\(p\) probability kernels.  Its integral over the \(p-1\) difference
variables and then the shifted label variable is one.  Tonelli's
theorem and successive integration of \(s_p,\ldots,s_2\), using
\eqref{l:finite-entropy:inverse}, therefore give
\[
 \int g_\eta^p\,d\rho\,dy\le B^{p-1};
\]
the last variable contributes \(\nu(\mathbb R^d)\le1\).

For a continuous compactly supported test function \(v(A,y)\),
integration against \(g_\eta\) converges to
\(\int v(A,As)\,d\rho(A)d\nu(s)\) as \(\eta\downarrow0\).
Shrinking support and bounded convergence prove this for the finite
measure \(\nu\).  H\"older bounds the limiting functional on
\(L^{p/(p-1)}(\rho(dA)\,dy)\) by \(B^{(p-1)/p}\).
Density of continuous compactly supported tests and \(L^p\) duality
give a representing density with that norm.  Equality on those tests
identifies it with the exact image measure, so a nonnegative version
may be used.  This proves \eqref{l:finite-entropy:projection-moment}.
The auxiliary kernels have been removed before any message is formed.
\end{proof}

\subsection{Information in one finite message}

\begin{lemma}[Regular block information]\label{l:finite-entropy:block}
There is an absolute \(c_*>0\) such that the following holds for
\(k=\lfloor c_*d\rfloor\).  Let \(W\) have any law \(P\) regular
on one cell of \(Q_{\mathrm{ho}}\), let \(A\) be an independent
\(k\)-by-\(d\) standard Gaussian matrix, and put \(Y=As(W)\).
Suppose \(U\) takes at most \(m_d\) values and is conditionally
independent of \(W\) given \((A,Y)\).  If
\(\log m_d=o(d^2)\), then
\begin{equation}\label{l:finite-entropy:block-information}
 I(W;U)\le2d
\end{equation}
for all sufficiently large \(d\), uniformly over these laws and
channels.  The threshold may depend on the alphabet sequence, while
\(c_*\) and the coefficient \(2\) are absolute.
\end{lemma}

\begin{proof}
For each positive-probability letter let
\(p_u=\Pr(U=u)\) and
\(\ell_u(w)=\Pr(U=u\mid W=w)\).  Its posterior density relative
to \(P\) is \(\ell_u/p_u\).  Set
\[
 B_u=\{w:\ell_u(w)/p_u>e^d\},\qquad \lambda=e^{-d}.
\]
Since this density integrates to one, \(P(B_u)\le\lambda\).
Outside \(W\in B_U\), the log posterior density is at most \(d\).
We will show that \(\Pr(W\in B_U)\) is exponentially small in
\(d\), and then bound its contribution using the finite alphabet.

Use \(z,R\) from Lemma~\ref{l:finite-entropy:ball}.
The exact joint law of \((A,Y)\) has a density \(D(A,y)\) by
Lemma~\ref{l:finite-entropy:projection} with \(a=0\).
The graph image of \(\lambda^{-1}P|_{B_u}\) has mass at most one
and satisfies the ball hypothesis with multiplier \(e^d\), namely
with \(a=1\).  Let \(g_u(A,y)\) be its exact projection density
from Lemma~\ref{l:finite-entropy:projection}.  Under the original
experiment,
\[
 b_u(A,Y):=\Pr(W\in B_u\mid A,Y)
 =\frac{\lambda g_u(A,Y)}{D(A,Y)}
\]
almost surely.  This is the Radon--Nikodym ratio of the restricted
joint measure to the original one; points where \(D=0\) have zero
original probability.

To control this ratio, choose an absolute
\(C_2>2\sqrt{2\pi e}\) and define
\[
 G=\{D(A,Y)\ge(C_2R)^{-k}\}.
\]
For every fixed \(s\) in the support, \(A(s-z)\) is centered
Gaussian with covariance at most \(R^2I_k\).  Markov's inequality
with exponential parameter \(1/4\) gives
\[
 \Pr\{\|Y-Az\|>2\sqrt kR\}\le e^{-k}2^{k/2}.
\]
The label ball of radius \(2\sqrt kR\) has volume at most
\((2\sqrt{2\pi e}R)^k\).  The probability of points inside it
where \(D<(C_2R)^{-k}\), after integrating over \(A\), is at
most \((2\sqrt{2\pi e}/C_2)^k\).  Thus
\begin{equation}\label{l:finite-entropy:good-probability}
 \Pr(G^c)\le2e^{-c_2k}
\end{equation}
for some absolute \(c_2>0\).

Let \(p=\lfloor\beta d/2\rfloor\).  On \(G\), the density
denominator is bounded below, so
\begin{equation}\label{l:finite-entropy:restricted-moment}
 \begin{split}
 \mathbb E[\mathbf1_G b_u(A,Y)^p]
 &\le \lambda^p(C_2R)^{k(p-1)}
              \int g_u(A,y)^p\,d\rho(A)\,dy\\
 &\le \exp\{-dp+C_3k(p-1)\},
 \end{split}
\end{equation}
where
\[
 C_3=\log C_2+(1+C_1)/\gamma+\log2
\]
is valid because \(k\ge1\).  Choose the fixed \(c_*>0\) so that
\[
 c_*\le\gamma/2,\qquad C_3c_*\le1/4.
\]
Then the projection lemma applies and the last moment is at most
\(e^{-3dp/4}\) for all sufficiently large \(d\).

Given \((A,Y)\), the conditional independence assumption lets us
average the choice of \(U\) separately from the event \(W\in B_u\).
Consequently
\[
 \Pr(W\in B_U,G)\le\mathbb E[\mathbf1_G\max_u b_u(A,Y)].
\]
The \(p\)th power of the maximum is at most the sum of the \(p\)th
powers.  Since \(\log m_d\le dp/4\) eventually, H\"older and
\eqref{l:finite-entropy:restricted-moment} bound this probability by
\(e^{-d/2}\).  Also \(k\ge c_*d/2\) eventually, so
\eqref{l:finite-entropy:good-probability} yields
\[
 q_0:=\Pr(W\in B_U)\le e^{-c_3d}
\]
after decreasing an absolute \(c_3>0\) and increasing the dimension
threshold.

For \(q_u=\Pr(U=u,W\in B_u)\le p_u\), the contribution of the
event \(W\in B_U\) to
\(\mathbb E\log(\ell_U(W)/p_U)\) is at most
\[
 \sum_uq_u\log(1/p_u)
 \le\sum_uq_u\log(1/q_u)
 \le q_0\log(m_d/q_0).
\]
The last bound is the entropy bound for at most \(m_d\) letters
after normalization by \(q_0\); its value is zero if \(q_0=0\).
The exponential bound on \(q_0\) and \(\log m_d=o(d^2)\) make
this \(o(1)\).  Outside the event, the log posterior density is at
most \(d\).  Hence \(I(W;U)\le d+o(1)\le2d\), as claimed.
\end{proof}

\subsection{Iteration in the finite-state model}

Use the finite-state model of Section~\ref{l:model}, now
with \(M(d)=o(d^2)\).  Suppose the original learner has angular
success at least \(2/3\) for every signal at \(0<\epsilon\le1/10\).
For the cube experiment, first fix the data-independent shared choice
of rules with average success greater than \(3/5\), which is possible
by averaging the original success \(2/3\).  The derivative of the graph
map contains the identity in its first \(n\) rows, so
Lemma~\ref{l:prior-ac} and then Lemma~\ref{l:borel} supply Borel
representatives for this fixed choice without changing its cube law.
Now implement the finite-output transition kernels by independent
uniforms indexed by state and layer, and sample the output law
independently for each terminal state and index.  Include the initial
state among these remaining choices; its conditional law given the
fixed shared seed is independent of the signal and rows by the joint
independence in Section~\ref{l:model}.  Averaging these choices fixes
a deterministic program with cube success at least
\(3/5\).  This reduction is only for the cube prior.

Use \(L_*=\lceil T/k\rceil\) blocks, padding the final block with
ignored independent rows.  Pregenerate rows after any early stop.
An active block message records its continuing state, or its terminal
state and within-block stopping offset, including immediate stopping.
Use a dummy message on all later blocks.  Thus the alphabet per block
has size at most
\begin{equation}\label{l:finite-entropy:alphabet}
 m_d=(k+2)2^M+1,\qquad \log m_d=o(d^2).
\end{equation}
The block number is already known, so no factor depending on \(T\)
is needed.  The messages determine the terminal state, global stopping
index, and fixed output.

Retain the messages in an analysis history.  Immediately after each
message, append the entire regularization label from
Lemma~\ref{l:finite-entropy:regularization}.  Initially the law is
\(\mu\), which is regular.  On every positive-probability history the
current law is regular on its indicated cell.  That history uses only
\(W\), past rows, and the fixed rules; future independent rows remain
independent of it and of \(W\).  Given the entering state, the next
message is a channel of the fresh block data.  Hence
Lemma~\ref{l:finite-entropy:block} applies conditionally; a dummy
message has zero information.

Let \(K_i\) be the expected conditional divergence to the uniform
law on the current cell after \(i\) messages and regularizations, and
let \(j_i\) be the expected current depth.  Both start at zero.
Conditioning a current density \(f\) on a message \(u\) multiplies
it by \(\ell_u/p_u\), where \(0\le\ell_u\le1\) and \(p_u>0\)
on that branch.  Taking logarithms and averaging shows that the
expected divergence to the old cell grows by exactly the conditional
mutual information of the message.  Finite entropy is preserved on
each branch: its positive log part grows by at most \(\log(1/p_u)\),
and its expectation under the new law is finite because that law has
density at most \(1/p_u\) relative to the old one.  Its negative
part has the universal integrable bound used above.
The average increase is at most \(2d\) by the block lemma.
Applying the regularization budget conditionally then gives
\begin{equation}\label{l:finite-entropy:iteration}
 K_{i+1}+bnj_{i+1}\le K_i+bnj_i+2d+2/e.
\end{equation}
This also proves inductively that the next finite-entropy hypothesis
holds.

Let \(\mathcal H\) be the full history, including every complement
decision. The bound just proved ensures finite expected divergence
and depth. Histories are countable, so averaging
\eqref{l:finite-entropy:cell-divergence} for their conditional laws
and using the defining Radon--Nikodym identity for mutual information
gives
\begin{equation}\label{l:finite-entropy:history-information}
 \begin{split}
 I(W;\mathcal H)
 &=K_{L_*}+n(\log2)j_{L_*}\\
 &\le\max\{1,\log2/b\}(2d+2/e)L_*
 \le C_4dL_*
 \end{split}
\end{equation}
for an absolute \(C_4\).  This calculation includes the information
in all auxiliary labels; it does not assume a bound on their individual
Shannon entropies.

The deterministic output is known from \(\mathcal H\).
Lemma~\ref{ld:cube-geometry} bounds the angular success probability
of any fixed output by \((\sqrt{2\pi e}\,\epsilon)^n\).
Under the product law of
\(W\) and \(\mathcal H\), the same bound holds; under the joint
law success is at least \(3/5\).  Equation~\eqref{l:binary} gives
the exact endpoint
\begin{equation}\label{l:finite-entropy:accuracy}
 I(W;\mathcal H)
 \ge\frac35n\log\frac1{\sqrt{2\pi e}\,\epsilon}-\log2.
\end{equation}
For \(L=\log(1/\epsilon)\ge\log10\),
\[
 L-\log\sqrt{2\pi e}
 \ge\left(1-\frac{\log\sqrt{2\pi e}}{\log10}\right)L,
\]
with a positive absolute coefficient.  Thus, for sufficiently large
\(d\), \eqref{l:finite-entropy:accuracy} is at least \(cdL\),
uniformly in the accuracy.  Combining it with
\eqref{l:finite-entropy:history-information} gives
\(\lceil T/k\rceil\ge c'L\).  When \(L\ge2/c'\),
\(T\ge k(c'L-1)\ge c''dL\), because \(k\ge c_*d/2\)
eventually.  When \(T=0\), the zero upper bound and positive lower
bound already contradict success.

For the bounded remaining range of \(L\), the original learner
satisfies \(T>d/100\).  Take a uniform full-sphere signal \(S\) and
grant the learner all \(T\le d/100\) pregenerated rows.  Let \(V\)
be their span and \(p_V=P_VS\).  Conditional on the rows and \(p_V\),
the residual is uniform on the sphere of radius
\(\sqrt{1-\|p_V\|^2}\) in \(V^\perp\).  This follows by writing
\(S\) as a normalized Gaussian: the direction of its orthogonal
Gaussian component is independent of both component lengths and of
the component in \(V\), hence also of the row-space projection after
normalization.  Isotropy gives
\[
 \mathbb E\|p_V\|^2\le T/d,\qquad
 \Pr\{\|p_V\|>1/2\}\le4T/d.
\]
Conditional also on the learner's independent randomness, its unit
output is fixed.  The inner product of that output with the residual
has mean zero and variance at most \(1/(d-T)\).  On
\(\|p_V\|\le1/2\), angular success requires this inner product to
be at least \(0.4\), since \(\cos(\epsilon)\ge0.9\).
Chebyshev's inequality therefore bounds success by
\begin{equation}\label{l:finite-entropy:linear}
 \frac{4T}{d}+\frac1{0.4^2(d-T)}<\frac23
\end{equation}
for sufficiently large \(d\), a contradiction.  Hence the bounded
range also has \(T=\Omega(dL)\).  All lower-bound constants are
absolute; the eventual dimension threshold may depend on the rate
\(M(d)/d^2\to0\).
